\documentclass[10pt,a4paper]{amsart}
\usepackage[margin=19mm]{geometry}
\usepackage{amsmath,amssymb,mathtools}
\usepackage[scr=rsfs,cal=euler]{mathalfa}
\usepackage{xfrac}
\usepackage[unicode,hidelinks,bookmarks=false]{hyperref}
\newcommand{\ie}{i.e.\ }
\newcommand{\cf}{cf.\ }
\newcommand{\resp}{respectively\ }
\newcommand{\Subsection}{\subsection}
\providecommand{\nobreakdash}{\nobreak\hbox{-}}
\setlength{\emergencystretch}{2em}
\multlinegap0pt
\usepackage{tikz}
\usetikzlibrary{cd}
\usepackage{amssymb}
\usepackage{mathtools}
\usepackage{dsfont}
\usepackage{xcolor}
\usepackage[shortlabels]{enumitem}
\newtheorem{lemma}{Lemma}
\newcommand{\C}{\mathbb{C}}
\newcommand{\Las}{\mathsf{L}}
\newcommand{\R}{\mathbf{R}}
\newcommand{\id}{\mathsf{id}}
\newcommand{\nto}{\mathrel{\mkern3mu\vcenter{\hbox{$\scriptscriptstyle|\:$}}\mkern-12mu{\to}}}
\title{Local categories: a new framework for partiality}
\author{Marcello Lanfranchi, Jean-Simon Pacaud Lemay}
\date{Source: arXiv:2512.03371v2 (CC BY 4.0)}
\begin{document}
\maketitle

\noindent\emph{Source and licence.} Local categories: a new framework for partiality. Version arXiv:2512.03371v2 (CC BY 4.0), distributed by arXiv under the Creative Commons Attribution 4.0 International licence, http://creativecommons.org/licenses/by/4.0/, as recorded in the arXiv licence metadata for this exact version and shown on the abstract page https://arxiv.org/abs/2512.03371v2. Full licence notice: \texttt{licenses/arxiv-2512.03371-CC-BY-4.0.txt}.\par\smallskip

\noindent\textit{Editorial scope.} Counted unit: the article's lemma lemma:split-local-category-2 (source0.tex lines 1769--1772). The article's complete original proof of it is delivered with it (source0.tex lines 1773--1876, one \texttt{proof} environment of 104 lines). No mathematics is written, repaired or re-derived: the statement and the proof are the article's own lines, in the article's own notation, and no retained line is substituted or re-flowed.  Retained uncounted as the source-local context the proof needs: lines 1760--1763.  Nothing was omitted from any retained span, so the package carries no omission record.  No retained line contains a citation, so the wrapper carries no bibliography and no bibliography entry.  The retained lines contain the article's own diagram code, which the wrapper typesets with the standard tikz packages; no image file is delivered and the package declares no asset dependency.  Editorial formatting only, in addition: an AMS \texttt{amsart} preamble that restates the article's own declarations and macros used by the retained lines, namely the environments \texttt{lemma} on separate counters, together with the standard packages those lines need.  The retained environments print in this document as Lemma 1 in order of appearance, whereas the article prints the same items with the numbering of its own full document.  The complete native source of the article as distributed in the arXiv e-print for this exact version is bundled unchanged as \texttt{sources/190262/source0.tex}.
\begin{lemma}
\label{lemma:monic-iso}
In any category, given two composable morphisms $f\colon A\to B$ and $g\colon B\to C$, if $fg$ is an isomorphism and $g$ is monic, then both $f$ and $g$ are isomorphisms.
\end{lemma}

\begin{lemma}
\label{lemma:split-local-category-2}
Let $\C$ be a local category. An object $U$ of $\C$ splits if and only if the restriction idempotent $(U,\eta_U)\colon A\nto A$ splits in $\R[\C]$, where $A=\Las U$. In particular, $\C$ is a split local category if and only if $\R[\C]$ is a split restriction category.
\end{lemma}

\begin{proof}
Assume that $U$ splits in $\C$. Thus, $U$ is isomorphic to a total object $E=\Las E$ via an isomorphism $s\colon E\to U$ and its inverse $r\colon U\to E$. Consider the morphisms $(U,r)\colon A\nto E$ and $(E,s\eta_U)\colon E\nto A$ of $\R[\C]$, where $A=\Las U$ and $\eta_U\colon U\to A$. Let us prove that $(U,r)$ and $(E,s\eta_U)$ define a splitting for $(U,\eta_U)\colon A\nto A$ in $\R[\C]$. For starters, let us compute the composition $(E,s\eta_U)(U,r)$:
\begin{equation*}
% https://q.uiver.app/#q=WzAsOCxbMCwwLCJFIl0sWzQsMCwiQSJdLFs2LDAsIkUiXSxbNSwxLCJVIl0sWzMsMSwiVSJdLFsyLDIsIkUiXSxbMywzLCJFIl0sWzQsMiwiVSJdLFszLDEsIlxcZXRhX1UiXSxbMywyLCJyIiwyXSxbNywzLCIiLDIseyJsZXZlbCI6Miwic3R5bGUiOnsiaGVhZCI6eyJuYW1lIjoibm9uZSJ9fX1dLFs0LDEsIlxcZXRhX1UiLDJdLFs3LDQsIiIsMSx7ImxldmVsIjoyLCJzdHlsZSI6eyJoZWFkIjp7Im5hbWUiOiJub25lIn19fV0sWzUsNCwicyIsMl0sWzYsNywicyIsMl0sWzYsNSwiIiwxLHsibGV2ZWwiOjIsInN0eWxlIjp7ImhlYWQiOnsibmFtZSI6Im5vbmUifX19XSxbNSwwLCIiLDEseyJsZXZlbCI6Miwic3R5bGUiOnsiaGVhZCI6eyJuYW1lIjoibm9uZSJ9fX1dLFs3LDEsIiIsMSx7InN0eWxlIjp7Im5hbWUiOiJjb3JuZXIifX1dLFs2LDQsIiIsMSx7InN0eWxlIjp7Im5hbWUiOiJjb3JuZXIifX1dXQ==
\begin{tikzcd}
E &&&& A && E \\
&&& U && U \\
&& E && U \\
&&& E
\arrow["{\eta_U}"', from=2-4, to=1-5]
\arrow["{\eta_U}", from=2-6, to=1-5]
\arrow["r"', from=2-6, to=1-7]
\arrow[equals, from=3-3, to=1-1]
\arrow["s"', from=3-3, to=2-4]
\arrow["\lrcorner"{anchor=center, pos=0.125, rotate=135}, draw=none, from=3-5, to=1-5]
\arrow[equals, from=3-5, to=2-4]
\arrow[equals, from=3-5, to=2-6]
\arrow["\lrcorner"{anchor=center, pos=0.125, rotate=135}, draw=none, from=4-4, to=2-4]
\arrow[equals, from=4-4, to=3-3]
\arrow["s"', from=4-4, to=3-5]
\end{tikzcd}
\end{equation*}
However, since $sr=\id_E$, we conclude that $(E,s\eta_U)(U,r)=\id_E$ in $\R[\C]$. We now compute the composition $(U,r)(E,s\eta_U)$ as follows:
\begin{equation*}
% https://q.uiver.app/#q=WzAsNixbMCwwLCJBIl0sWzIsMCwiRSJdLFs0LDAsIkEiXSxbMSwxLCJVIl0sWzMsMSwiRSJdLFsyLDIsIlUiXSxbMywwLCJcXGV0YV9VIl0sWzMsMSwiciIsMl0sWzQsMiwic1xcZXRhX1UiLDJdLFs0LDEsIiIsMCx7ImxldmVsIjoyLCJzdHlsZSI6eyJoZWFkIjp7Im5hbWUiOiJub25lIn19fV0sWzUsNCwiciIsMl0sWzUsMywiIiwyLHsibGV2ZWwiOjIsInN0eWxlIjp7ImhlYWQiOnsibmFtZSI6Im5vbmUifX19XSxbNSwxLCIiLDIseyJzdHlsZSI6eyJuYW1lIjoiY29ybmVyIn19XV0=
\begin{tikzcd}
A && E && A \\
& U && E \\
&& U
\arrow["{\eta_U}", from=2-2, to=1-1]
\arrow["r"', from=2-2, to=1-3]
\arrow[equals, from=2-4, to=1-3]
\arrow["{s\eta_U}"', from=2-4, to=1-5]
\arrow["\lrcorner"{anchor=center, pos=0.125, rotate=135}, draw=none, from=3-3, to=1-3]
\arrow[equals, from=3-3, to=2-2]
\arrow["r"', from=3-3, to=2-4]
\end{tikzcd}
\end{equation*}
However, since $rs=\id_U$, we conclude that $(U,r)(E,s\eta_U)=(U,\eta_U)$. Thus, $((E,s\eta_U),(U,r))$ forms a splitting for $(U,\eta_U)$ in $\R[\C]$.\\
Conversely, assume that $(U,\eta_U)$ splits in $\R[\C]$. Thus, there are two morphisms $(U_1,r')\colon A\nto E$ and $(U_2,s')\colon E\to A$ such that $(U_1,r')(U_2,s')=(U,\eta_U)$ and $(U_2,s')(U_1,r')=\id_E$. From these two equations, we deduce the existence of four morphisms of $\C$ as follows:
\begin{align*}
&p_1\colon U\to U_1             &&p_2\colon U\to U_2\\
&q_1\colon E\to U_1             &&q_2\colon E\to U_2
\end{align*}
satisfying the following conditions
\begin{equation*}
% https://q.uiver.app/#q=WzAsMyxbMCwwLCJVIl0sWzEsMCwiVV8xIl0sWzEsMSwiQSJdLFsxLDIsIlxcZXRhXzEiXSxbMCwxLCJwXzEiXSxbMCwyLCJcXGV0YV9VIiwyXV0=
\begin{tikzcd}
U & {U_1} \\
& A
\arrow["{p_1}", from=1-1, to=1-2]
\arrow["{\eta_U}"', from=1-1, to=2-2]
\arrow["{\eta_1}", from=1-2, to=2-2]
\end{tikzcd}\qquad
% https://q.uiver.app/#q=WzAsMyxbMCwwLCJVIl0sWzEsMCwiVV8yIl0sWzEsMSwiQSJdLFsxLDIsInMnIl0sWzAsMSwicF8yIl0sWzAsMiwiXFxldGFfVSIsMl1d
\begin{tikzcd}
U & {U_2} \\
& A
\arrow["{p_2}", from=1-1, to=1-2]
\arrow["{\eta_U}"', from=1-1, to=2-2]
\arrow["{s'}", from=1-2, to=2-2]
\end{tikzcd}\qquad
% https://q.uiver.app/#q=WzAsMyxbMCwwLCJFIl0sWzEsMCwiVV8xIl0sWzEsMSwiRSJdLFsxLDIsInInIl0sWzAsMSwicV8xIl0sWzAsMiwiIiwyLHsibGV2ZWwiOjIsInN0eWxlIjp7ImhlYWQiOnsibmFtZSI6Im5vbmUifX19XV0=
\begin{tikzcd}
E & {U_1} \\
& E
\arrow["{q_1}", from=1-1, to=1-2]
\arrow[equals, from=1-1, to=2-2]
\arrow["{r'}", from=1-2, to=2-2]
\end{tikzcd}\qquad
% https://q.uiver.app/#q=WzAsMyxbMCwwLCJFIl0sWzEsMCwiVV8yIl0sWzEsMSwiRSJdLFsxLDIsIlxcZXRhXzIiXSxbMCwxLCJxXzIiXSxbMCwyLCIiLDIseyJsZXZlbCI6Miwic3R5bGUiOnsiaGVhZCI6eyJuYW1lIjoibm9uZSJ9fX1dXQ==
\begin{tikzcd}
E & {U_2} \\
& E
\arrow["{q_2}", from=1-1, to=1-2]
\arrow[equals, from=1-1, to=2-2]
\arrow["{\eta_2}", from=1-2, to=2-2]
\end{tikzcd}
\end{equation*}
where, as a shorthand, we denoted by $\eta_1$ and $\eta_2$, $\eta_{U_1}$ and $\eta_{U_2}$, respectively. Furthermore, the following two diagrams are pullbacks in $\C$:
\begin{equation*}
% https://q.uiver.app/#q=WzAsNCxbMCwwLCJVIl0sWzEsMCwiVV8yIl0sWzAsMSwiVV8xIl0sWzEsMSwiRSJdLFswLDIsInBfMSIsMl0sWzAsMSwicF8yIl0sWzEsMywiXFxldGFfMiJdLFsyLDMsInInIiwyXSxbMCwzLCIiLDEseyJzdHlsZSI6eyJuYW1lIjoiY29ybmVyIn19XV0=
\begin{tikzcd}
U & {U_2} \\
{U_1} & E
\arrow["{p_2}", from=1-1, to=1-2]
\arrow["{p_1}"', from=1-1, to=2-1]
\arrow["\lrcorner"{anchor=center, pos=0.125}, draw=none, from=1-1, to=2-2]
\arrow["{\eta_2}", from=1-2, to=2-2]
\arrow["{r'}"', from=2-1, to=2-2]
\end{tikzcd}\qquad
% https://q.uiver.app/#q=WzAsNCxbMCwwLCJFIl0sWzEsMCwiVV8xIl0sWzAsMSwiVV8yIl0sWzEsMSwiQSJdLFswLDEsInFfMSJdLFswLDIsInFfMiIsMl0sWzIsMywicyciLDJdLFsxLDMsIlxcZXRhXzEiXSxbMCwzLCIiLDEseyJzdHlsZSI6eyJuYW1lIjoiY29ybmVyIn19XV0=
\begin{tikzcd}
E & {U_1} \\
{U_2} & A
\arrow["{q_1}", from=1-1, to=1-2]
\arrow["{q_2}"', from=1-1, to=2-1]
\arrow["\lrcorner"{anchor=center, pos=0.125}, draw=none, from=1-1, to=2-2]
\arrow["{\eta_1}", from=1-2, to=2-2]
\arrow["{s'}"', from=2-1, to=2-2]
\end{tikzcd}
\end{equation*}
From $q_2\eta_2=\id_E$, that $\eta_2$ is monic, and by Lemma~\ref{lemma:monic-iso}, we deduce that both $q_2$ and $\eta_2$ are isomorphisms in $\C$. However, since $p_1$ is a pullback of $\eta_2$ and $\eta_2$ is invertible, $p_1$ is also an isomorphism, since isomorphisms are stable under pullbacks. From the equation $p_2s'=\eta_U$, since $\eta_U$ is monic, so is $p_2$. However, since $p_1r'=p_2\eta_2$, that $p_1$ is invertible and that $\eta_2$ is monic, we conclude that $r'=p_1^{-1}p_2\eta_2$ is also monic. Finally, using that $q_1r'=\id_E$, that $r'$ is monic, and by Lemma~\ref{lemma:monic-iso}, $r'$ is an isomorphism in $\C$. Thus, since both $p_1\colon U\to U_1$ and $r'\colon U_1\to E$ are isomorphisms, $U$ is isomorphic to $E$ in $\C$. However, since $E$ is an object of $\R[\C]$, $E$ is total in $\C$.
\end{proof}

\end{document}
